\section{Lecture 10: Support Vector Machines (Conceptual)}

\subsection{Motivation}

In previous lectures, we studied:
\begin{itemize}
    \item Linear and logistic regression
    \item Feature maps and kernel methods
\end{itemize}

\medskip

\noindent
These methods aim to separate data using linear decision boundaries.

\medskip

\noindent
However, a key question remains:

\medskip

\noindent
\textit{Among all separating hyperplanes, which one should we choose?}

\medskip

\noindent
\textbf{Guiding idea.}  
Choose the decision boundary that is most robust to perturbations.

\subsection{Margin Maximization}

Consider binary classification with labels $y_i \in \{-1, +1\}$.

\medskip

\noindent
A linear classifier is:
\[
f(x) = w^\top x + b.
\]

\medskip

\noindent
\textbf{Decision boundary.}
\[
w^\top x + b = 0.
\]

\medskip

\noindent
\textbf{Margin.}
\begin{itemize}
    \item Distance between the decision boundary and the closest data points
\end{itemize}

\medskip

\noindent
\textbf{Geometric margin.}
\[
\gamma = \min_i \frac{y_i (w^\top x_i + b)}{\|w\|}.
\]

\medskip

\noindent
\textbf{Optimization problem.}
\[
\max_{w,b} \; \gamma.
\]

\medskip

\noindent
\textbf{Equivalent formulation.}
\[
\min_{w,b} \; \|w\|^2 \quad \text{subject to} \quad y_i(w^\top x_i + b) \geq 1.
\]

\medskip

\noindent
\textbf{Interpretation.}
\begin{itemize}
    \item Maximizing margin improves robustness
    \item Only closest points matter
\end{itemize}

\medskip

\noindent
\textbf{Support vectors.}
\begin{itemize}
    \item Points that lie on the margin boundary
    \item Determine the classifier
\end{itemize}

\medskip

\noindent
\textbf{Insight.}  
SVM focuses on the most informative data points.

\subsection{Hard and Soft Margins}

\textbf{Hard-margin SVM.}
\begin{itemize}
    \item Assumes perfectly separable data
    \item Constraints:
    \[
    y_i(w^\top x_i + b) \geq 1
    \]
\end{itemize}

\medskip

\noindent
\textbf{Limitation.}
\begin{itemize}
    \item Sensitive to noise
    \item Often infeasible in real data
\end{itemize}

\medskip

\noindent
\textbf{Soft-margin SVM.}

Introduce slack variables $\xi_i \geq 0$:
\[
y_i(w^\top x_i + b) \geq 1 - \xi_i.
\]

\medskip

\noindent
\textbf{Objective.}
\[
\min_{w,b,\xi} \; \|w\|^2 + C \sum_i \xi_i.
\]

\medskip

\noindent
\textbf{Interpretation.}
\begin{itemize}
    \item Allows violations of margin constraints
    \item Balances margin size and classification error
\end{itemize}

\medskip

\noindent
\textbf{Connection to regularization.}
\begin{itemize}
    \item $\|w\|^2$ acts as regularizer
    \item $C$ controls tradeoff between fit and complexity
\end{itemize}

\medskip

\noindent
\textbf{Insight.}  
Soft margins extend SVM to noisy, real-world data.

\subsection{Dual Formulation (Conceptual)}

Rather than solving directly for $w$, we can express the solution in terms of training data.

\medskip

\noindent
\textbf{Key idea.}
\begin{itemize}
    \item The optimal $w$ can be written as:
    \[
    w = \sum_{i=1}^n \alpha_i y_i x_i.
    \]
\end{itemize}

\medskip

\noindent
\textbf{Interpretation.}
\begin{itemize}
    \item Only support vectors have $\alpha_i > 0$
    \item Solution depends on a subset of data points
\end{itemize}

\medskip

\noindent
\textbf{Inner products.}
\[
w^\top x = \sum_i \alpha_i y_i (x_i^\top x).
\]

\medskip

\noindent
\textbf{Insight.}  
The classifier depends only on inner products between data points.

\subsection{Kernel Methods Revisited}

From the dual perspective, predictions depend on:
\[
x_i^\top x.
\]

\medskip

\noindent
We replace this with a kernel:
\[
k(x_i, x) = \langle \phi(x_i), \phi(x) \rangle.
\]

\medskip

\noindent
\textbf{Kernel SVM.}
\[
f(x) = \sum_i \alpha_i y_i k(x_i, x) + b.
\]

\medskip

\noindent
\textbf{Effect.}
\begin{itemize}
    \item Linear classifier in feature space
    \item Nonlinear classifier in input space
\end{itemize}

\medskip

\noindent
\textbf{Key observation.}
\begin{itemize}
    \item Kernel trick enables high-dimensional modeling efficiently
\end{itemize}

\medskip

\noindent
\textbf{Insight.}  
SVM combines margin maximization with kernel-based representation.

\subsection{A Unifying Perspective}

Support vector machines unify multiple ideas:

\begin{itemize}
    \item \textbf{Geometry:} margin maximization
    \item \textbf{Optimization:} constrained quadratic problem
    \item \textbf{Regularization:} control of model complexity
    \item \textbf{Representation:} kernel feature spaces
\end{itemize}

\medskip

\noindent
\textbf{Core idea.}  
A good classifier is one that separates data with maximum robustness.

\medskip

\noindent
\textbf{Key insight.}  
The combination of margin and representation leads to powerful nonlinear models.

\subsection{Outlook}

In this lecture, we introduced:

\begin{itemize}
    \item margin maximization,
    \item hard and soft margin SVM,
    \item dual formulation (conceptually),
    \item kernel methods in SVM.
\end{itemize}

\medskip

\noindent
These ideas illustrate how geometry, optimization, and representation interact.

\medskip

\noindent
In the next lecture, we study model evaluation and selection, focusing on how to assess and compare models in practice.

\medskip

\noindent
\textbf{Guiding principle.}  
Effective learning balances data fit, robustness, and representation.